\hypertarget{python-3.13-free-threaded-build-gil-removal-internals}{%
\chapter{Python 3.13: Free-Threaded Build \& GIL Removal
Internals}\label{python-3.13-free-threaded-build-gil-removal-internals}}

\hypertarget{the-free-threaded-cpython-paradigm-shift}{%
\subsection{20.1 The Free-Threaded CPython Paradigm
Shift}\label{the-free-threaded-cpython-paradigm-shift}}

Introduced provisionally in Python 3.13 (PEP 703), the free-threaded
build of CPython allows executing bytecode in parallel across multiple
OS threads without the constraint of the Global Interpreter Lock (GIL).
To maintain thread safety and prevent race conditions on shared objects,
CPython's runtime engine has been fundamentally rewritten, moving away
from single-threaded assumptions to scalable multi-threaded internals.

\begin{center}\rule{0.5\linewidth}{0.5pt}\end{center}

\hypertarget{biased-reference-counting-brc}{%
\subsection{20.2 Biased Reference Counting
(BRC)}\label{biased-reference-counting-brc}}

Reference counting is the primary mechanism for memory management in
CPython. In a standard build, incrementing/decrementing reference counts
is a simple, non-atomic operation:
\passthrough{\lstinline!op->ob\_refcnt++!}. In a thread-parallel
environment, standard operations would cause race conditions. Using
standard atomic operations (\passthrough{\lstinline!atomic\_add!} /
\passthrough{\lstinline!atomic\_sub!}) globally would destroy
performance due to cache coherency traffic (bus locks) across CPU cores.

To solve this, PEP 703 introduces \textbf{Biased Reference Counting
(BRC)}. BRC divides an object's reference count state based on thread
ownership.

\hypertarget{biased-object-headers}{%
\subsubsection{1. Biased Object Headers}\label{biased-object-headers}}

An object is ``biased'' toward the thread that allocated it (the owner
thread). The owner thread uses fast, non-atomic CPU instructions to
modify the reference count. Non-owner threads (foreign threads) must use
atomic operations on a separate ``shared'' reference count field.

\hypertarget{c-level-struct-representation-of-reference-layout}{%
\subsubsection{2. C-level Struct representation of reference
layout}\label{c-level-struct-representation-of-reference-layout}}

In the free-threaded CPython header definition
(\passthrough{\lstinline!object.h!} under no-GIL configurations), the
object header is redefined:

\begin{lstlisting}[language=C]
typedef struct _object {
    uintptr_t ob_ref_local;     /* Local reference count accessed only by owner thread */
    uintptr_t ob_ref_shared;    /* Shared reference count accessed atomically by foreign threads */
    PyTypeObject *ob_type;      /* Pointer to object type descriptor */
} PyObject;
\end{lstlisting}

\hypertarget{brc-reference-counting-algorithm}{%
\subsubsection{3. BRC Reference Counting
Algorithm}\label{brc-reference-counting-algorithm}}

When a reference is modified, CPython checks the current thread ID
against the owner thread ID encoded inside the object's header:
\[\text{Is Owner} = (\text{Current Thread ID} == \text{ob\_ref\_local} \ \& \ \text{THREAD\_ID\_MASK})\]

\begin{lstlisting}
                  [Reference Count Change Request]
                                 |
                     Is Current Thread the Owner?
                               /   \
                             Yes    No
                             /       \
     [Non-Atomic Local Increment]   [Atomic Shared Increment]
      (Fast path: CPU registers)     (Slow path: CAS / Bus Lock)
\end{lstlisting}

If the owner thread detects that the local reference count drops to
zero, it merges the shared reference count. If the combined total
reference count reaches zero, the object is safely deallocated.

\begin{center}\rule{0.5\linewidth}{0.5pt}\end{center}

\hypertarget{thread-safe-allocation-via-mimalloc}{%
\subsection{20.3 Thread-Safe Allocation via
mimalloc}\label{thread-safe-allocation-via-mimalloc}}

CPython's classic allocator (\passthrough{\lstinline!pymalloc!}) is
optimized for small objects but is single-threaded and not thread-safe.
In the free-threaded build, \passthrough{\lstinline!pymalloc!} is
replaced by \textbf{mimalloc}, a highly efficient, thread-safe,
metadata-compact memory allocator developed by Microsoft. *
\textbf{Thread-Local Pools}: Mimalloc organizes memory in thread-local
heaps, allowing allocation to proceed without acquiring global locks in
most cases. * \textbf{Thread-Safe Free Lists}: When a foreign thread
deallocates an object, it does not write directly back to the owner
thread's main free list. Instead, it writes to a thread-safe ``atomic
free list'' associated with the target page, preventing cross-thread
allocation bottlenecks.

\begin{center}\rule{0.5\linewidth}{0.5pt}\end{center}

\hypertarget{hazard-pointers-and-deferred-reference-counting}{%
\subsection{20.4 Hazard Pointers and Deferred Reference
Counting}\label{hazard-pointers-and-deferred-reference-counting}}

Some global or long-lived objects (such as interned strings, code
constants, or modules) are read frequently by many threads but rarely
modified. To bypass BRC overhead on these objects: 1. \textbf{Immortal
Objects}: CPython marks these objects as immortal by setting a special
bitmask in their reference count field. The interpreter completely skips
reference counting modifications for immortal objects. 2. \textbf{Hazard
Pointers}: When traversing lock-free data structures (like internal
dictionaries or type caches), threads use \textbf{Hazard Pointers} to
declare they are accessing an object. This ensures the object will not
be freed by the GC or another thread while in use, avoiding the need to
perform expensive atomic reference count updates for transient reads.

\begin{center}\rule{0.5\linewidth}{0.5pt}\end{center}
